\chapter{Caste, Democracy and Caste Associations}

\section{Caste and Democracy}

\subsection{Answering 'Sanskritization vs Westernization'}

\begin{ExampleBox}
\begin{itemize}
\item A model answer to 'Differentiate between Sanskritization and Westernization' (10 marks):
\item \textbf{Introduce} --- M. N. Srinivas introduced both concepts to explain social change. \textbf{Sanskritization} = imitation of the cultural lifestyle of the upper castes by low castes/tribals; \textbf{Westernization} = a process of institutional, ideological and technological change during 150 years of British rule.
\item \textbf{Contrast} --- Sanskritization promotes a \textbf{sacred outlook} (copying Brahminical lifestyle) and \textbf{mobility within caste}; Westernization promotes a \textbf{secular outlook} (egalitarianism, meritocracy) and \textbf{mobility outside caste}. Sanskritization was availed by the \textbf{low castes}; Westernization by the \textbf{upper castes} (primary westernization).
\item \textbf{Dynamic relation} --- Westernization can be a prerequisite to Sanskritization; they are also opposed.
\item \textbf{Critique} --- Yogendra Singh: Sanskritization is orthogenetic change, Westernization is heterogenetic.
\item The lesson: \textbf{address the demand} --- a 4th/5th-attempt student who dumps the whole Coorg case study and the three dimensions of Westernization (institutional/ideological/technological) fails the 'difference' question, while a first-attempt student who sticks to the contrast scores. Competition is about \textbf{addressing the demand}, not knowing more.
\end{itemize}
\end{ExampleBox}

\subsection{Further Critiques: Dube, Oommen, and Dumont}

\begin{CriticalPoint}
\begin{itemize}
\item \textbf{S. C. Dube} (field study in \textbf{Shamirpet, Telangana}) --- against the dominant-caste concept: 'where is the dominant caste --- village, district, region, state or national?' Srinivas did not clarify; one village's dominant caste is not the next village's. Moreover, it is \textbf{not the caste but individuals/families} that are dominant (e.g. Mulayam Singh Yadav's family in UP/Bihar).
\item \textbf{T. K. Oommen} (recap) --- \textbf{multiple power structures}: 'power reservoir' (a community that is a pool of power --- e.g. the \textbf{Meena} among STs; farmers as a reservoir) and 'power exercisers' (the few who actually exercise that power --- e.g. the \textbf{Tikaits} in the farmers' protest).
\end{itemize}
\end{CriticalPoint}

\begin{CriticalPoint}
\begin{itemize}
\item \textbf{Critique of Louis Dumont}:
\item 1. His understanding of caste is based on a \textbf{static, book view} --- he did not see the \textbf{evolution of caste} (unlike Ghurye); caste is no longer rigid.
\item 2. The \textbf{purity/impurity} category projects the \textbf{Brahminical view} of Hindu society, ignoring the \textbf{subjectivities} of purity (Dalits also claim themselves pure).
\item 3. He \textbf{ignored the protest movements} in Indian history that questioned the ideology of caste division, projecting Hindu religious values as a universal category (as if there is no conflict against those values).
\item 4. \textbf{McKim Marriott} --- \emph{Homo Hierarchicus} is a \textbf{'speculative sketch of a pair of models shaped by the textual ideology of social science'} --- Dumont's models (homo hierarchicus/aequalis/major/minor) are ideal types projected as pure/absolute types.
\item 5. \textbf{Nicholas Dirks} (\emph{Castes of Mind}) --- challenged Dumont's 'power is inferior to status' as absurd: \textbf{purohits were only advisors to kings} (kings employed them); there is no \textbf{diametrically opposite relation} between status and power. (Analogy: is the Prime Minister inferior to the President? It depends on the scale --- head of government vs head of state.) Purohits served kings --- because \textbf{power matters}. Brahmin rulers were rare in ancient India (only the Satavahanas claimed Brahminical status).
\item 6. \textbf{Veena Das} --- in villages the relation between castes is \textbf{not always hierarchical}: there is also \textbf{coordination and cooperation} --- a relation of \textbf{reciprocity} (e.g. the jajmani relation is patron--client, not purely hierarchical).
\end{itemize}
\end{CriticalPoint}

\subsection{Caste and Democracy: The Classical Discourse}

\begin{KeyConcept}
\begin{itemize}
\item \textbf{Ambedkar} believed \textbf{Hinduism and democracy are opposite} --- democracy is about equality, while Hinduism is established on the inequality of caste.
\item Initially it was believed that \textbf{caste and democracy are two opposite social categories} --- democracy as an open system that would weaken or abolish caste.
\item But caste became a reality in India \textbf{only with the British} --- the \textbf{caste census} became an instrument that \textbf{created caste} (classifying people into high and low, giving the hierarchy an official stamp).
\end{itemize}
\end{KeyConcept}

\begin{KeyConcept}
\begin{itemize}
\item The \textbf{classical sociologists} (Ghurye, Srinivas, Dumont, Beteille) initiated the \textbf{discourse on caste and democracy} by examining the \textbf{dynamic relation (static and fluid)} between them.
\item This analysis helped the \textbf{modern political sociologists} (the \textbf{Rudolphs}) to frame a \textbf{hypothesis}: '\textbf{with more democratization, you get more caste associations}' --- i.e. more horizontal solidarities.
\end{itemize}
\end{KeyConcept}

\subsection{Ghurye: Caste Patriotism}

\begin{KeyConcept}
\begin{itemize}
\item \textbf{Ghurye} was the first to discuss the \textbf{politicization of caste} (the phrase was coined by \textbf{Rajni Kothari}).
\item Traditionally \textbf{caste had integrated India}; with the British and the \textbf{caste census}, caste assumed a \textbf{new (political) avatar}.
\item In \textbf{1932} the non-Brahmins mobilized against the Brahmins in South India (the census revealed Brahmin dominance in wealth, education and representation). On the surface this seemed to weaken the hierarchy, but it actually \textbf{generated a new collective sentiment --- caste solidarity} --- which Ghurye termed \textbf{caste patriotism}.
\end{itemize}
\end{KeyConcept}

\subsection{Srinivas: From Vertical to Horizontal Solidarity}

\begin{KeyConcept}
\begin{itemize}
\item \textbf{Srinivas elaborated Ghurye's thesis}: '\textbf{far from disappearing with modernization, caste was now experiencing horizontal consolidation}'.
\item \textbf{Before independence} there was \textbf{vertical solidarity} --- the interdependence of upper and lower castes (the jajmani relation; in Coorg and Rampura, Brahmins depended on the lower castes and vice versa).
\item \textbf{After independence}, vertical solidarity weakened and was replaced by \textbf{horizontal solidarity}: castes formed their \textbf{own associations} (Jat Sabha, All-India Thakur Sabha, Brahman Sabha, Reddy association).
\item How did this happen? The \textbf{Constitution gave everyone the right to form associations} (a fundamental right); \textbf{democracy runs on competition}, so these consolidated caste associations compete with each other --- democracy turned out to be, in effect, \textbf{competition between caste associations} rather than between free individuals.
\item This is a \textbf{change in the structure, not a change of the structure} --- caste remains; only the form of solidarity has changed, from vertical (across castes) to horizontal (within the caste).
\end{itemize}
\end{KeyConcept}

\subsection{Dumont and Beteille: Substantialization and the New Caste}

\begin{KeyConcept}
\textbf{Dumont} analysed the same process and gave it a new name: \textbf{'caste did not disappear with economic and political change, but its logic definitely got altered'} --- earlier caste was a \textbf{structure}, now it has become a \textbf{substance} (impenetrable universal blocks). This is the \textbf{'substantialization of caste'} (as more modernization, more substantialization).
\end{KeyConcept}

\begin{KeyConcept}
\begin{itemize}
\item \textbf{Beteille's} one-line answer on caste and democracy: \textbf{before democracy Shripuram had a highly undifferentiated caste structure (caste was class); after modernization, the old caste was replaced by a NEW caste} --- the title of his book \textbf{'Caste, Old and New'}.
\item The \textbf{new caste} = \textbf{status groups} --- English-speaking professionals for whom \textbf{ritual knowledge is no longer important}; what matters is party etiquette, the way you speak English and dress. This \textbf{achieved-status group has emerged as the new caste}.
\item With democracy and modernization, the undifferentiated Shripuram \textbf{started differentiating}: \textbf{class and power, once embedded in caste, gained relative autonomy and became disembedded} --- Shripuram was waiting for a class structure, but got the new caste instead.
\item \textbf{Beteille's broader contribution}: he was the \textbf{first to introduce Weberian sociology into India} (before him, Indian sociology had Orientalism --- Ghurye; Marxism --- A. R. Desai; and structural functionalism --- Srinivas, but no Max Weber).
\end{itemize}
\end{KeyConcept}

\begin{QuickRevision}
\begin{itemize}
\item Answer strategy: address the demand (Sanskritization = sacred/mobility within caste/low castes; Westernization = secular/mobility outside/upper castes; dynamic relation; Yogendra Singh's critique).
\item S. C. Dube: 'where is the dominant caste?' + dominant individuals/families; T. K. Oommen: multiple power structures (reservoir vs exercisers; Meena, Tikaits).
\item Critique of Dumont: static book view; purity/impurity = Brahminical view (ignores Dalit subjectivity); ignores protest movements; Marriott (speculative models); Nicholas Dirks (power not inferior to status; purohits served kings); Veena Das (reciprocity, not hierarchy, in villages).
\item Caste \& democracy: Ambedkar (Hinduism vs democracy opposite); caste census created caste; Ghurye (politicization, caste patriotism, 1932 non-Brahmin mobilization).
\item Srinivas: vertical solidarity (pre-independence) $\rightarrow$ horizontal solidarity (caste associations post-independence); democracy = competition between associations; change in structure, not of structure.
\item Dumont: substantialization (structure $\rightarrow$ substance); Beteille: old caste $\rightarrow$ new caste (English-speaking professionals/status groups); class/power disembedded from caste.
\end{itemize}
\end{QuickRevision}

\section{Caste Associations and the Evolution of Caste Politics}

\subsection{What is a Caste Association?}

\begin{KeyConcept}
\begin{itemize}
\item \textbf{'Tradition gave us caste, modernity gave us democracy, and Indian modernity gave us caste association.'}
\item A \textbf{caste association is a formal organization of similar jatis having similar interests} (e.g. the Reddy association, Vanniyar association, Vadama association, Maratha sabha, Brahman sabha, Thakur sabha, Jat sabha/mahasabha). It acts as a \textbf{pressure group}, especially at election time.
\item It is a \textbf{bureaucracy} --- a formal, registered body (registered with the Government of India, receiving funds) with \textbf{hierarchy and division of labour} (a Jat leader, his associates, volunteers, IT/WhatsApp workers). It is also a \textbf{'party'} in Weber's sense (people with common interests), though not a political party.
\end{itemize}
\end{KeyConcept}

\subsection{Functions of Caste Associations}

\begin{KeyConcept}
\begin{itemize}
\item The functions of a caste association:
\item 1. \textbf{Welfare of its jati and its members} --- social, economic, political and educational: helping poor members, rebuilding a member's house after a fire, running hostels and coaching, and \textbf{ensuring representation in politics} (an association loses power if its members are not represented in the legislature/panchayat/administration). It works like a \textbf{fraternity} --- e.g. coming together to defend a member facing charges.
\item 2. \textbf{Justice to the community} (closely tied to welfare).
\item 3. \textbf{Improving the social status of the community} --- the census as an instrument of social mobility.
\end{itemize}
\end{KeyConcept}

\begin{ExampleBox}
\textbf{The Vanniyar (Palli) case}: the Palli of northern Tamil Nadu, originally agricultural labourers, did not want to be classified as Shudra. Through their association they claimed \textbf{Kshatriya} status, asserting descent from the \textbf{Pallava dynasty / Agni Kula}; a court case ruled in their favour, and the census recorded them as Kshatriya. So the \textbf{census both creates caste and helps its mobility} --- an \textbf{unanticipated consequence} (Merton).
\end{ExampleBox}

\subsection{Origin and Phase 1: Urban Welfare Associations}

\begin{KeyConcept}
\begin{itemize}
\item Caste associations emerged with \textbf{British colonialism}: modern technology (travel, meetings, pamphlets, print) made it possible to form all-India associations. \textbf{Westernization consolidated caste identity} --- and democracy promoted it (the right to form associations). At the national level they appeared from the \textbf{1890s} onwards (locally from the 1870s).
\item \textbf{Phase 1 (late 19th century, 1890s--1930s) --- URBAN caste associations}: they ran \textbf{cooperative societies} (housing, hospitals, educational institutions, hostels), and worked to \textbf{improve their caste's status} (claiming Kshatriya status, demanding temple entry, getting their names and histories recorded in the census --- the \textbf{Jati Purans}). They performed \textbf{both ritual and secular functions}, but this was largely restricted to urban pockets.
\item (Caste mobility is not only Hindu: among Muslims the lower \textbf{Arzals} emulating the upper \textbf{Ashrafs} is called \textbf{Ashrafization}.)
\end{itemize}
\end{KeyConcept}

\subsection{Phase 2: Poona Pact and Pressure Groups (1930s--70s)}

\begin{KeyConcept}
\begin{itemize}
\item \textbf{Phase 2 (1930s--1970s)}: the \textbf{Poona Pact (1932)} clearly separated Indians into \textbf{Savarna and Avarna} --- earlier the separation was ritual, now it became \textbf{legal} too, which \textbf{strengthened Dalit identity}.
\item Ambedkar offered to lead the Dalits; Gandhi's leadership was rejected by many; and when Dalits became \textbf{westernized}, a hierarchy emerged within them --- \textbf{Dalit elites} (e.g. Mayawati) who are a distinct faction within the Dalits.
\item With the \textbf{growth of representational politics, castes began to act as pressure groups} --- the Dalit movement, the Depressed Class Association, the Dalit Panthers.
\item In \textbf{1967} Congress faced its first defeat (in eight states); \textbf{factions/lobbies} of backward castes emerged within Congress, and the \textbf{locally dominant castes, empowered by land reforms, entered national and regional politics}. In Rajni Kothari's words, \textbf{'caste gets politicized'} --- and 'Congress as a system' began to be disrupted.
\end{itemize}
\end{KeyConcept}

\subsection{Phase 3: Dalit Politics and BAMCEF (1970s)}

\begin{KeyConcept}
\begin{itemize}
\item \textbf{Phase 3 (1970s)}: rise of \textbf{regional parties and decline of Congress}. The agrarian (dominant) castes formed their own parties or powerful factions within Congress.
\item From \textbf{1977--78}, new caste associations came from the \textbf{low castes (Dalits)}: \textbf{BAMCEF} (Backward and Minority Communities Employees Federation), founded by \textbf{Kanshi Ram}, paved the way for the \textbf{Bahujan Samaj Party (BSP)} --- which, in the wake of the \textbf{Mandal Commission report}, became a catalyst for the emergence of the \textbf{OBCs} as a significant political force.
\item Dalit consciousness was now \textbf{on democratic lines}, organised around instruments provided by democracy (the political party). Politics came to be led by \textbf{cross-caste platforms} using the language of \textbf{assertion and resistance} for greater inclusion of oppressed groups.
\item The \textbf{demand of caste associations changed} --- from \textbf{collective welfare to political activism} (a social change from pre- to post-independence).
\end{itemize}
\end{KeyConcept}

\subsection{Phase 4: Mandalization of Politics (1980s--90s)}

\begin{KeyConcept}
\begin{itemize}
\item \textbf{Phase 4 (1980s--90s) --- 'Mandalization of politics'}: a \textbf{rebellion against upper-caste dominance at the national level}, expressed through the weakening of national parties (like Congress --- the 'Brahmin-Bania party') and the rise of regional parties.
\item When \textbf{OBC reservation was implemented (1991)}, the \textbf{anti-reservation agitation} had an \textbf{unanticipated consequence}: it united the \textbf{SC, ST and OBC} (the 'class unity' of the backward classes) --- but it also \textbf{revived the question of caste and politics and normalized caste as the mode of politics}.
\item A scholar noted that \textbf{caste groups, instead of crumbling, adapted themselves to the demands of parliamentary politics} --- producing a \textbf{'democracy of caste'}: competition between \textbf{caste groups}, not between individuals. ('We don't cast our vote; we vote our caste.')
\item The shift was from the \textbf{politics of ideology to the politics of representation}.
\end{itemize}
\end{KeyConcept}

\subsection{Rudolph \& Rudolph and Rajni Kothari}

\begin{KeyConcept}
\begin{itemize}
\item \textbf{Rudolph and Rudolph} (Lloyd and Susanne Rudolph, 1967): the \textbf{novelty of the caste association} was the \textbf{combination of the ascriptive dimension of caste with voluntary association} --- caste is ascriptive (membership by birth), but membership of the caste association is \textbf{achieved/voluntary} (annual fee, elections, democratic participation). This is the \textbf{coexistence of tradition and modernity}. Caste also serves as an \textbf{associational platform} (providing networks for organizations).
\item \textbf{Rajni Kothari} (post-1990s): it seems at the surface that our politics is 'caste-ridden', but the reality is the reverse --- what we are experiencing is the \textbf{'democratization of caste'}: caste has assumed a \textbf{non-ritual, purely political function}, and has been \textbf{democratized}. Caste and democracy are not opposite --- they form a new mix.
\end{itemize}
\end{KeyConcept}

\begin{CriticalPoint}
In conclusion, the evolution of caste associations shows a tendency of \textbf{group loyalties over merit and competence}, and \textbf{caste patriotism over public interest} --- both of which \textbf{run counter to liberal and democratic values}.
\end{CriticalPoint}

\begin{QuickRevision}
\begin{itemize}
\item Caste association = formal organization of similar jatis with similar interests (pressure group; a Weberian 'party'; a bureaucracy).
\item Functions: welfare (social/economic/political/educational) + justice + improving social status (census = mobility; Vanniyar/Palli $\rightarrow$ Kshatriya via court + census; unanticipated consequence).
\item Origin: colonial modern technology; westernization consolidated caste; national-level from 1890s.
\item Phase 1 (1890s--1930s): urban associations, cooperative societies (housing/hospital/education/hostels), status-raising (Kshatriya claims, temple entry, census), ritual + secular.
\item Phase 2 (1930s--70s): Poona Pact (legal Savarna/Avarna split); Dalit elites; castes as pressure groups; 1967 Congress defeat; dominant castes (land reforms) enter national politics; 'caste gets politicized' (Kothari).
\item Phase 3 (1970s): regional parties; BAMCEF (Kanshi Ram) $\rightarrow$ BSP; Mandal $\rightarrow$ OBC force; Dalit consciousness on democratic lines; cross-caste platforms (assertion/resistance); welfare $\rightarrow$ political activism.
\item Phase 4 (1980s--90s): Mandalization of politics (rebellion vs upper-caste dominance; Congress decline); anti-reservation agitation united SC/ST/OBC but normalized caste as mode of politics; democracy of caste.
\item Rudolph \& Rudolph: ascriptive caste + voluntary association; Kothari: democratization of caste (not caste-ridden politics).
\item Conclusion: group loyalties > merit; caste patriotism > public interest; runs counter to democracy.
\end{itemize}
\end{QuickRevision}

